\SourceLeaf{386}%
% A544-P0386-D01
[I.] Art.\ $17$. Die erste Umformung des letzten $S$ in dem Ausdrucke für $\Delta S$ erhält man, wenn man das betreffende Flächenstück in solche drei Theile zerlegt, dass jenes $S$ in
\[
\begin{aligned}
&S\bigl(\tfrac12m+\tfrac12-\tfrac12i,\ (\tfrac12+\tfrac12i)m,\\
&\qquad(\tfrac12+\tfrac12i)m+i,\ \tfrac12m+\tfrac12+\tfrac12i_{(+)}\bigr)\\
&-S\bigl(\tfrac12m,\ \tfrac12m+\tfrac12+\tfrac12i_{(+)},\ \tfrac12m+\tfrac12-\tfrac12i\bigr)\\
&-S\bigl((\tfrac12+\tfrac12i)m,\ \tfrac12m,\ \tfrac12m+\tfrac12-\tfrac12i\bigr)
\end{aligned}
\]
übergeht, und wenn man dann die Ganzepunkte in dem ersten Flächentheile mit Hülfe des Satzes in Art.\ $10$ auszählt und ferner berücksichtigt, dass in dem zweiten Flächentheile sich kein Ganzepunkt befindet.

% A544-P0386-D02
Die zweite Umformung erhält man, wenn man die den Eckpunkten des dritten Flächentheils entsprechenden Grössen mit $i$ multiplicirt und um die ganze Zahl $(1-i)\frac{m-1}{2}$ vermehrt, endlich die dritte Umformung, wenn man mit der zuletzt entstandenen Figur nach Vorschrift des Art.\ $16$ diejenige vergleicht, die gegen jene die Ortsverschiedenheit $\frac{-1}{1+i}$ hat.

% A544-P0386-D03
[I.] Art.\ $[18.]$ Eine Erläuterung zum ersten Schema findet man in dem später niedergeschriebenen hier mit [II.] bezeichneten Bruchstücke Art.\ $1$ bis $6$.

% A544-P0386-D04
[I.] Art.\ $[18.]\ (2.)$ Die geometrische Deutung ergiebt mit Zuhülfenahme der beiden Systeme von Ganzepunkten
\[
[-2iQ-\tfrac12,\tfrac12,iQ]=[-2iQ+\tfrac12,\tfrac12,iQ]=\mathrm{IV}^{*}
\]
und
\[
-[-2iQ+\tfrac12i,-\tfrac12i,iQ]=-[-2iQ-\tfrac12i,-\tfrac12i,iQ]=\mathrm{X}^{*}
\]
die Gleichungen
\[
\begin{aligned}
\mathrm{IV}-\mathrm{IV}^{*}+[-2iQ,\tfrac12i,-\tfrac12]
&=-\mathrm{X}+\mathrm{X}^{*}+[-iQ,\tfrac12i,-\tfrac12]\\
\mathrm{XIII}+\mathrm{IV}^{*}&=[-2iQ,1,iQ]=-I(-2iQ)+I(-iQ)\\
\mathrm{V}+\mathrm{X}^{*}&=-[-2iQ,-i,iQ]=R(-2iQ)-R(-iQ)
\end{aligned}
\]
wenn allgemein $Rx$ und $Ix$ die grössten Ganzen des reellen Theils und des Coëfficienten des imaginären Theils von $x$ bedeuten.

% A544-P0386-D05
$[-2iQ,\tfrac12i,-\tfrac12]$ ist aber die Anzahl der Ganzepunkte in dem Quadrate, dessen Mittelpunkt sich in $-2iQ-\tfrac14+\tfrac14i$ befindet und zwischen dessen Endpunkten die Ortsunterschiede $\tfrac12$ und $\tfrac12i$ Statt haben.

% A544-P0386-D06
$[-iQ,\tfrac12i,-\tfrac12]$ oder $[-2iQ,-\tfrac12i,\tfrac12]$ ist die Anzahl der Ganzepunkte in einem gleichen Quadrate mit dem Mittelpunkte $-2iQ+\tfrac14-\tfrac14i$.

% A544-P0386-D07
[I.] Art.\ $[18.]\ (3.)$ Mit Zuhülfenahme der Ganzepunkte $[0,\tfrac12i,-iQ]=-[-\tfrac12i,-\tfrac12i,-iQ]=\mathrm{I}^{*}$ erhält man

\SourceLeaf{387}%
% A544-P0387-D01
\[
\begin{aligned}
\mathrm{VII}-\mathrm{XII}&=\mathrm{I}-\mathrm{I}^{*}+[-iQ,\tfrac12,\tfrac12i]-[0,\tfrac12,\tfrac12i]\\
&=\mathrm{I}-\mathrm{I}^{*}+[-2iQ,\tfrac12,\tfrac12i]\\
\mathrm{II}-\mathrm{I}^{*}&=[0,-i,-iQ]=-R0+R(-iQ)\\
\mathrm{IX}-\mathrm{I}&=[0,-1,-iQ]=I0-I(-iQ)\\
\mathrm{VI}-\mathrm{XII}&=-[\tfrac12i,-1,-iQ]=-I(\tfrac12i)+I(-iQ+\tfrac12i)\\
\mathrm{VIII}-\mathrm{XIV}&=-[\tfrac12,-1-i,\tfrac12im_{(-)}]=\frac{a-1}{2}+\frac b2
\end{aligned}
\]

% A544-P0387-D02
[II.] Art.\ $10$. Es ist
\[
\begin{aligned}
\Upsilon P'&\equiv-\Upsilon(-iP'+\tfrac12im),\\
\Upsilon P''&\equiv-1-\Upsilon(P''-\tfrac12im),\\
\Upsilon P'''&\equiv1+\Upsilon(-iP''')\ (\text{mod. }4)
\end{aligned}
\]
und $-iP'+\tfrac12im$, $P''-\tfrac12im$ sind die um $\tfrac12i$ vermehrten Ganzepunkte resp.\ in I, VI.

% A544-P0387-D03
[III.] Art.\ $6$. Die in der Note angegebenen Regeln für die Bestimmung des $\mathrm{Dec.}\frac Mm$ habe ich der vorliegenden Abhandlung aus einem andern Orte der Handschriften beigefügt. Die erste dieser beiden Regeln, die wie leicht zu sehen mit der zweiten übereinstimmt, folgt aus der des Art.\ $6$, weil
\[
\begin{gathered}
k\equiv f\pdot i^{-n}\ (\text{mod. }m),\quad n=\Theta\frac{2fM}{m},\quad p=mm',\\
\left[\frac{2km'}p\right]^2\equiv\Theta\frac{2km'}p\ (\text{mod. }2+2i)
\end{gathered}
\]
ist.

% A544-P0387-D04
[III.] Art.\ $8$ enthält in der Handschrift ein Beispiel zu Art.\ $7$, nemlich die Bestimmung des Decidenten von $-1+2i$ für den Modulus $-11+4i$.

% A544-P0387-D05
[III.] Art.\ $10$. In Bezug auf die Bemerkung ‘anders auszudrücken’ kann man Art.\ $3$ des folgenden Bruchstücks [IV] vergleichen.

% A544-P0387-D06
[IV.] Die Art.\ $1$. $2$. $4$ enthalten in der Handschrift ausser dem hier Abgedruckten noch die Anwendung auf die beiden Beispiele für $m=5+8i$, $M=9+4i$ und für $m=9+4i$, $M=5+8i$.

% A544-P0387-D07
[V.] Art.\ $[7.]$ Es bezeichnet hier $\mathrm{Dec.}\frac mM$ wie in Art.\ $1$ des vorhergehenden Bruchstücks [IV] den Werth von
\[
\Sigma(-1)^X\Sigma(-1)^Y\mathrm{Int.}\,p
\]
worin die Summation über alle ganze Zahlen $X$ und $Y$ auszudehnen ist, für welche die zugehörigen $\xi$ und $\eta$ innerhalb der Grenzen $0$ und $\tfrac12$ liegen.

% A544-P0387-D08
Die Formeln für den Decidenten in Art.\ $7$ und $15$ sind nach der Angabe des Textes auf zwei be%
\SourceLeaf{388}%
% A544-P0388-D01
sondern Wegen gefunden, um aber diese Erläuterungen nicht zu sehr auszudehnen, werden sie hier aus einer gemeinsamen Quelle abgeleitet.

% A544-P0388-D02
Indem $X$ irgend einen bestimmten ganzzahligen Werth annimmt, sei $Y^{*}$ das kleinere, $Y^{**}$ das grössere der beiden $Y$, welche den Grenzwerthen von $\xi$, $\eta$ entsprechen. Die zu $Y^{*}$ und $Y^{**}$ zugehörigen Werthe von $p$ seien $p^{*}$ und $p^{**}$, die ebenso wie $Y^{*}$ und $Y^{**}$ einander nicht gleich werden können, weil die Summe $\Sigma$ sich nicht über die Grenzwerthe von $\xi$ und $\eta$ erstreckt.

% A544-P0388-D03
Führt man auf dieselbe Weise wie in den beiden vorhergehenden Aufsätzen [III] und [IV] die Summation über alle bei demselben $X$ Statt habenden Werthe von $Y$ aus, setzt dabei für die Anzahl der zwischen $Y'$ und $Y''$ liegenden ungeraden Zahlen $[\tfrac12Y''-\tfrac12]-[\tfrac12Y'-\tfrac12]$ und fügt die Intensoren, die sich auf die Grenzen $\xi=0$ und $=\tfrac12$ beziehen, zwei Mal aber mit entgegengesetzten Zeichen hinzu, so erhält man für $\Sigma(-1)^Y\mathrm{Int.}\,p$ den aus sieben Theilen bestehenden Ausdruck
\[
\begin{gathered}
-\Sigma[-\mathrm{Int.}(p-\mu\omega i)+\mathrm{Int.}(p+\mu\omega i)]\\
\text{worin alle }p\text{ aufzunehmen, für welche }[Y]\text{ gerade,}\\
x\text{ oder }y\text{ ganz, incl. }\xi=0\text{ und }\tfrac12,\quad
\text{excl. }\eta=0\text{ und }\tfrac12
\end{gathered}
\]
\[
\left.\begin{aligned}
&-\mathrm{Int.}(p^{*}-\mu\omega i)\text{ wenn }[Y^{*}]\text{ gerade}\\
&+\mathrm{Int.}(p^{**}+\mu\omega i)\text{ wenn }[Y^{**}]\text{ gerade}
\end{aligned}\right\}\ \xi=0\text{ oder }\tfrac12,\quad0<\eta<\tfrac12
\]
\[
\left.\begin{aligned}
&-\mathrm{Int.}(p^{*}+\mu\omega i)\text{ wenn }[Y^{*}]\text{ gerade}\\
&+\mathrm{Int.}(p^{**}-\mu\omega i)\text{ wenn }[Y^{**}]\text{ gerade}
\end{aligned}\right\}\ 0<\xi<\tfrac12,\quad\eta=0\text{ oder }\tfrac12
\]
\[
\left.\begin{aligned}
&-\mathrm{Int.}(p^{*}+\mu\omega i)\text{ wenn }[Y^{*}]\text{ gerade}\\
&+\mathrm{Int.}(p^{**}-\mu\omega i)\text{ wenn }[Y^{**}-\omega]\text{ gerade}
\end{aligned}\right\}\ \xi=0\text{ oder }\tfrac12,\quad\eta=0\text{ oder }\tfrac12
\]
welcher mit $(-1)^X$ multiplicirt und über alle ganzzahligen $X$ summirt den Decidenten $\frac mM$ ergiebt.

% A544-P0388-D04
Aus dem ersten Theil des Ausdrucks entsteht auf diese Weise von den nach Art.\ $2$ Vorschrift I gebildeten $\varepsilon$
\[
\begin{aligned}
&\Sigma\varepsilon\text{ wo }X\text{ ganz, }[Y]\text{ gerade}\\
&-\Sigma\varepsilon\text{ wo ausserdem }y\text{ ganz, }[x]\text{ gerade}
\end{aligned}
\]
Für die folgenden Theile kann
\[
\left.\begin{aligned}
&-(B)\mathrm{Int.}(p-B\mu\omega i)\text{ wenn }\xi=0\\
&+(B)\mathrm{Int.}(p+B\mu\omega i)\text{ wenn }\xi=\tfrac12
\end{aligned}\right\}\ X\text{ ganz, }[Y]\text{ gerade, }0<\eta<\tfrac12
\]
\[
\left.\begin{aligned}
&-(A)\mathrm{Int.}(p+A\mu\omega i)\text{ wenn }\eta=0\\
&+(A)\mathrm{Int.}(p-A\mu\omega i)\text{ wenn }\eta=\tfrac12
\end{aligned}\right\}\ X\text{ ganz, }[Y]\text{ gerade, }0<\xi<\tfrac12
\]
\[
\begin{gathered}
-\tfrac12[(A')+(B')]\mathrm{Int.}(p+A'\mu\omega i)\\
\text{wenn }X\text{ ganz, }[Y+A'\omega]\text{ gerade, }\xi\text{ und }\eta=0\text{ oder }\tfrac12
\end{gathered}
\]
gesetzt werden, worin $A'=+A$ oder $-A$ ist, wenn $\eta=0$ oder $\tfrac12$, $B'=+B$ oder $-B$, wenn $\xi=0$ oder $\tfrac12$, und worin z.\ B.\ $(A):+1$ oder $-1$ bezeichnet, jenachdem $A$ positiv oder negativ ist.

% A544-P0388-D05
Multiplicirt man mit $(-1)^X$, führt die Summation über $X$ aus, lässt dabei in diesen Ausdrücken u.\ zwar im
\[
\begin{aligned}
\text{ersten}\quad&p-B\mu\omega i,\ P-BD\omega i,\ \pi-AB\omega i-BB\omega,\ X,\ [Y]\\
&\text{bez. in }ip,\ iP,\ i\pi,\ -Y,\ [X]\\[3pt]
\text{zweiten}\quad&p+B\mu\omega i,\ P+BD\omega i,\ \pi+AB\omega i+BB\omega,\ X,\ [Y]\\
&\dotstwo\quad p,\ P,\ \pi\quad X,\ [Y]\\[3pt]
\text{dritten}\quad&p+A\mu\omega i,\ P+AD\omega i,\ \pi+AA\omega i+AB\omega,\ X,\ [Y]\\
&\dotstwo\quad p,\ P,\ \pi\quad X,\ [Y]\\[3pt]
\text{vierten}\quad&p-A\mu\omega i,\ P-AD\omega i,\ \pi-AA\omega i-AB\omega,\ X,\ [Y]\\
&\dotstwo\text{ in }im-ip,\ iM-iP,\ i-i\pi,\ Y-B,\ A-1-[X]
\end{aligned}
\]

\SourceLeaf{389}%
% A544-P0389-D01
übergehen und bezeichnet das aus dem fünften Ausdruck sich ergebende Resultat mit $Q_1$, so entsteht
\[
\begin{aligned}
&-\Sigma(-1)^Y(B)\mathrm{Int.}\,ip,\\
&\qquad\text{wo }Y\text{ ganz, }[X]\text{ gerade, }\eta=\omega,\ 0<\xi<\tfrac12\\
&+\Sigma(-1)^X(B)\mathrm{Int.}\,p,\\
&\qquad\text{wo }X\text{ ganz, }[Y]\text{ gerade, }\xi=\tfrac12+\omega,\ 0<\eta<\tfrac12\\
&-\Sigma(-1)^X(A)\mathrm{Int.}\,p,\\
&\qquad\text{wo }X\text{ ganz, }[Y]\text{ gerade, }\eta=\omega,\ 0<\xi<\tfrac12\\
&+\Sigma(-1)^Y(A)\mathrm{Int.}(im-ip),\\
&\qquad\text{wo }Y\text{ ganz, }[X]\text{ gerade, }\xi=\tfrac12+\omega,\ 0<\eta<\tfrac12\\
&+Q_1
\end{aligned}
\]

% A544-P0389-D02
Die Untersuchung der einzelnen Fälle lässt erkennen, dass unter der Voraussetzung $M\equiv1\ (\text{mod. }2+2i)$
\[
\begin{aligned}
Q_1={}&-\mathrm{Int.}\,\mu\omega i\text{ ist, wenn }M\text{ im }1.\text{ Quadranten liegt}\\
&-\mathrm{Int.}(\tfrac12mi+\mu\omega i),\quad\text{wenn }M\text{ im }2.\text{ Quadr. und }\frac{M-1}{2}\text{ gerade}\\
&+\mathrm{Int.}(\tfrac12mi-\mu\omega i),\quad\text{wenn }M\text{ im }4.\text{ Quadr. und }\frac{M-1}{2}\text{ gerade}
\end{aligned}
\]
indem man eine complexe Zahl gerade oder ungerade nennt, je nachdem sie durch $2$ theilbar ist oder nicht.

% A544-P0389-D03
Hiernach wird also bei Anwendung der in den Vorschriften II und IV bestimmten $\varepsilon$
\[
\begin{aligned}
\mathrm{Dec.}\frac mM={}&\mathrm{I},& &\Sigma\varepsilon,\ \text{wo }X\text{ ganz, }[Y]\text{ gerade}\\
&&&-4\Sigma\varepsilon,\ \text{wo noch }y\text{ ganz, }[x]\text{ gerade}\\
&\mathrm{II},&&-\Sigma\varepsilon\mathrm{Int.}\,ip,\ \text{wo }Y\text{ ganz, }[X]\text{ gerade}\\
&\mathrm{IV},&&+\Sigma\varepsilon\mathrm{Int.}\,p,\ \text{wo }X\text{ ganz, }[Y]\text{ gerade}\\
&\mathrm{II},&&+\Sigma\varepsilon\mathrm{Int.}\,p,\ \text{wo }X\text{ ganz, }[Y]\text{ gerade}\\
&\mathrm{IV},&&+\Sigma\varepsilon\mathrm{Int.}(im-ip),\ \text{wo }Y\text{ ganz, }[X]\text{ gerade}\\
&&&+Q_1
\end{aligned}
\]

% A544-P0389-D04
In einer andern Form erhält man den Ausdruck für den Decidenten, wenn man zuerst nach $X$ summirt und dabei die Anzahl der zwischen $X'$ und $X''$ liegenden ungeraden Zahlen durch $[\tfrac12X''+\tfrac12]-[\tfrac12X'+\tfrac12]$ darstellt, nemlich
\[
\begin{aligned}
\mathrm{Dec.}\frac mM={}&\mathrm{I},& &\Sigma\varepsilon,\ \text{wo }Y\text{ ganz, }[X]\text{ ungerade}\\
&&&-4\Sigma\varepsilon,\ \text{wo noch }y\text{ ganz, }[x]\text{ gerade}\\
&\mathrm{II},&&-\Sigma\varepsilon\mathrm{Int.}\,ip,\ \text{wo }X\text{ ganz, }[Y]\text{ gerade}\\
&\mathrm{IV},&&+\Sigma\varepsilon\mathrm{Int.}\,p,\ \text{wo }Y\text{ ganz, }[X]\text{ ungerade}\\
&\mathrm{II},&&+\Sigma\varepsilon\mathrm{Int.}\,p,\ \text{wo }Y\text{ ganz, }[X]\text{ ungerade}\\
&\mathrm{IV},&&+\Sigma\varepsilon\mathrm{Int.}(im-ip),\ \text{wo }X\text{ ganz, }[Y]\text{ ungerade}\\
&&&+Q_2
\end{aligned}
\]
\[
\begin{aligned}
Q_2={}&-\mathrm{Int.}(-\mu\omega),\quad\text{wenn }M\text{ im }2.\text{ Quadr.}\\
&+\mathrm{Int.}(\tfrac12m-\mu\omega),\quad\text{wenn }M\text{ im }1.\text{ Quadr. und }\frac{M-1}{2}\text{ ungerade}\\
&-\mathrm{Int.}(\tfrac12m+\mu\omega),\quad\text{wenn }M\text{ im }3.\text{ Quadr. und }\frac{M-1}{2}\text{ ungerade}
\end{aligned}
\]

\SourceLeaf{390}%
% A544-P0390-D01
Führt man die Summation nach $Y$ zuerst aus, wählt aber die zweite so eben angewandte Art der Bestimmung der Anzahl der zwischen zwei Werthen liegenden ungeraden Zahlen, so wird
\[
\begin{aligned}
\mathrm{Dec.}\frac mM={}&\mathrm{I},& &\Sigma\varepsilon,\ \text{wo }X\text{ ganz, }[Y]\text{ ungerade}\\
&&&-4\Sigma\varepsilon,\ \text{wo noch }y\text{ ganz, }[x]\text{ gerade}\\
&\mathrm{II},&&-\Sigma\varepsilon\mathrm{Int.}\,ip,\ \text{wo }Y\text{ ganz, }[X]\text{ ungerade}\\
&\mathrm{IV},&&+\Sigma\varepsilon\mathrm{Int.}\,p,\ \text{wo }X\text{ ganz, }[Y]\text{ ungerade}\\
&\mathrm{II},&&+\Sigma\varepsilon\mathrm{Int.}\,p,\ \text{wo }X\text{ ganz, }[Y]\text{ ungerade}\\
&\mathrm{IV},&&+\Sigma\varepsilon\mathrm{Int.}(im-ip),\ \text{wo }Y\text{ ganz, }[X]\text{ ungerade}\\
&&&+Q_3
\end{aligned}
\]
\[
\begin{aligned}
Q_3={}&-\mathrm{Int.}(-\mu\omega i),\quad\text{wenn }M\text{ im }3.\text{ Quadr.}\\
&-\mathrm{Int.}(\tfrac12mi+\mu\omega i),\quad\text{wenn }M\text{ im }2.\text{ Quadr. und }\frac{M-1}{2}\text{ ungerade}\\
&+\mathrm{Int.}(\tfrac12mi-\mu\omega i),\quad\text{wenn }M\text{ im }4.\text{ Quadr. und }\frac{M-1}{2}\text{ ungerade}
\end{aligned}
\]

% A544-P0390-D02
Summirt man zuerst noch $X$ und gebraucht dabei die erste Art der Darstellung der Anzahl der zwischen zwei Werthen liegenden ungeraden Zahlen, so erhält man die in [V.] Art.\ $15$ angegebene Form für den Decidenten, wo die Grösse $q$ auch durch folgende Gleichung definirt werden kann
\[
\begin{aligned}
q={}&-\mathrm{Int.}\,\mu\omega,\quad\text{wenn }M\text{ im }4.\text{ Quadr.}\\
&+\mathrm{Int.}(\tfrac12m-\mu\omega),\quad\text{wenn }M\text{ im }1.\text{ Quadr. und }\frac{M-1}{2}\text{ gerade}\\
&-\mathrm{Int.}(\tfrac12m+\mu\omega),\quad\text{wenn }M\text{ im }3.\text{ Quadr. und }\frac{M-1}{2}\text{ gerade}
\end{aligned}
\]

% A544-P0390-D03
Die Vereinigung dieser vier Ausdrücke für den Decidenten bildet das in [V.] Art.\ $7$ aufgestellte Resultat, weil $\mathrm{Int.}\,ip-\mathrm{Int.}\,p$ gleich $3$ wird für $[x]$ gerade $[y]$ ungerade, sonst aber gleich $1$, ferner $\mathrm{Int.}(im-ip)+\mathrm{Int.}\,p$ gleich $0$ für $[x]$ gerade $[y]$ gerade, in den übrigen Fällen aber gleich $4$.

% A544-P0390-D04
[V.] Art.\ $[7.]$ Die erste Tafel für das Beispiel giebt in der ersten Spalte die zu jedem ganzzahligen $P$ zugehörigen Werthe von $\frac{37\pdot P}{M}$ oder $37(\xi+\eta i)$, wenn $0<\xi<\tfrac12$, $0<\eta<\tfrac12$ ist, in der zweiten $\frac{37\pdot Pm}{M}$ oder $37\pdot p$, in der dritten die in $p$ enthaltene grösste ganze Zahl, in der vierten $\pm\mathrm{Int.}\,p$, wo das obere Zeichen gilt, wenn $P$ durch $1+i$ theilbar, das untere, wenn $P$ nicht durch $1+i$ theilbar ist.

% A544-P0390-D05
[V.] Art.\ $[9.]$. $[12.]$ Die verbesserte Bezeichnungsweise der $\theta$ ist nur bei der zweiten und dritten Classe Artt.\ $9$. $10$ angedeutet, aber auch auf die erste und vierte Artt.\ $8$. $11$ auszudehnen. Hiernach wird ein $\theta^\lambda=T+Ui$ denjenigen Index $\lambda,=0,1,2$ oder $3$ haben, für welchen die durch die Gleichungen

